
Dentro de los \textit{RTOS}, las tareas deben cumplir con una restricción temporal. Según la manera de evaluar el cumplimiento de dicha restricción, es decir, su exactitud, estas pueden ser: por un lado \textit{Hard} y por otro \textit{Soft} y \textit{Firm}.

En las \textit{Hard Tasks}, si no cumplen con la \textit{deadline} impuesta, provocan un fallo catastrófico en el peor de los casos, y en el mejor de ellos una degradación del sistema \cite{Lipari2015}. En las \textit{Soft Tasks}, se tolera que una fracción de las tareas puedan no llegar a cumplir la restricción temporal anteriormente impuesta \cite{Erickson2022}. En cambio, en las \textit{Firm Tasks}, por cada conjunto de tareas, al menos un número de ellas deben cumplir con la \textit{deadline} \cite{Behrouzian2020}.

Para la implementación de circuitos biohíbridos con interacciones realistas, es necesario hacer uso de \textit{RTOS} orientados a las \textit{Hard} o \textit{Firm tasks}, puesto que los \textit{Hybrots} tienen ventanas de tiempo del orden de milisegundos o menos, dado el tiempo transcurrido entre la adquisición del impulso y la interacción con la neurona en los experimentos biológicos \cite{Amaducci2019}, por lo que necesitan atender todas las tareas dentro de dicha restricción (o al menos un gran número de ellas) para poder realizar la medición y sincronización con dichas células. Además de esto, se han elegido sistemas código abierto y uso gratuito, que pueden merecer la pena a la hora de explotar las capacidades de la computación en tiempo real.
